% RESEARCH_DRAFT_SHA256: 5bebf3607395b6cc802648f535032418f2fc26d0ca6aaa468363d1ad266afbda
% Layout and source-reference wrappers expanded for the web renderer.
\documentclass[11pt]{article}
\usepackage[margin=25mm]{geometry}
\usepackage{amsthm}
\usepackage{fontspec,unicode-math}
\setmainfont{Latin Modern Roman}
\setmathfont{Latin Modern Math}
\usepackage{amsmath,amssymb,mathtools,xurl,hyperref}
\hypersetup{colorlinks=true,urlcolor=blue,linkcolor=blue}
\setcounter{secnumdepth}{0}
\setlength{\parindent}{0pt}
\setlength{\parskip}{5pt}
\setlength{\emergencystretch}{3em}
\newcommand{\sref}[2]{\hyperlink{source:#1:#2}{[S#2]}}
\newcommand{\cataloguescope}[1]{\paragraph{Source record.}#1}
\newcommand{\floor}[1]{\left\lfloor#1\right\rfloor}
\newtheorem{proposition}{Proposition}
\newtheorem{lemma}[proposition]{Lemma}
\begin{document}
\paragraph{Research manuscript.} OpenAI \texttt{gpt-6-astra}, reasoning \texttt{ultra}; 11 October 2026.
\paragraph{Preserved source body.} The complete catalogue mathematical text and sources follow with layout and source links expanded, including source qualifications and its historical preparation notes. New research begins at the section headed Research attempt.

% UnsolvedMath ID 20002120; source catalogue code AIM-INFRASTRUCTURE-0020
\section{A testable training-support-renewal framework for online instructional tools}\label{20002120}
{\small UnsolvedMath ID 20002120 · Computer Science}
\subsection{Definitions and mathematical statement}

\paragraph{Shared notation.}Standard number systems are $\mathbb N,\mathbb Z,\mathbb Q,\mathbb R,\mathbb C$. Specialized notation and index ranges follow the cited source and the definitions below.


PD means professional development: organized training intended to improve instructors' knowledge or practice. Training on an online tool concerns its use; ongoing support can include technical assistance, instructional resources, and opportunities to adapt teaching. A support intervention and its intended educational outcome must be distinguished.
\paragraph{Record scope and qualifications.}
The catalogue title is retained as metadata. Named certificates, contracts, bounds or frameworks in that title are not asserted as results of the original source extract.
The registry marks this wording as reconstructed and unverified. No specific training package, resource constraint, instructor population, or success criterion is given.

Catalogue screening limitation: Educational assessment, adoption, or study-design discussion; the supplied question does not state a mathematical target with defined assumptions and outcome.

The source snapshot labels this record partially solved; this is catalogue metadata rather than independent verification of a mathematical result or the current completion of a project.
The source text below is rendered literally where necessary; damaged formulas and composite extracts are not silently repaired into a new mathematical claim.
The supplied catalogue marks this formulation as an unverified reconstruction; missing assumptions are not certified by the source link.
\paragraph{Statement recorded in the supplied source.}
22. a) PD for instructors, on the use of online tools

b) What other support is needed? \href{http://aimath.org/pastworkshops/oerassessmentproblems.pdf}{[S2]}
\paragraph{Source status.}
The supplied catalogue labels this source record partially solved; that metadata was checked against the supplied snapshot on 2026-10-11. The label is preserved for provenance and is not an independent verification of present open status or mathematical correctness.
\subsection{Short English statement}
Identify professional development and additional support needed by instructors using online tools.
\subsection{Sources}
\begin{itemize}
\item[S1] Ulam AI and the UnsolvedMath Contributors, UnsolvedMath: A Curated Collection of Open Mathematics Problems, record 20002120 (AIM-INFRASTRUCTURE-0020). Catalogue curation and metadata: CC BY 4.0. The source extract is retained; definitions, English summaries and qualifications are editorial additions.
\url{https://huggingface.co/datasets/ulamai/UnsolvedMath}.
\item[S2] AIM Workshop Problem Lists, source linked by catalogue record AIM-INFRASTRUCTURE-0020.
\url{http://aimath.org/pastworkshops/oerassessmentproblems.pdf}.
\end{itemize}
\paragraph{Preparation and review.}
The original source and UnsolvedMath contributors retain attribution for the source extract. Codex (GPT-6; exact serving revision not exposed) assisted with definitions, English summaries, scope notes and formatting. No new human mathematical authorship or independently verified solution is asserted. Permanent source ID 20002120 and code AIM-INFRASTRUCTURE-0020 checked against the complete supplied registry and current repository. Distinct upstream records, including repeated questions and broad or damaged extracts, are retained for complete coverage of the source CS category. A repeated record is not a new independent conjecture. Complete mathematical-equivalence and current-literature review has not been performed. Preparation and final review follow contribution prompts C and D, with source gaps disclosed. Structural and rendering checks do not establish a new theorem or open status.
\label{end:20002120}

\section{Research attempt}
\subsection{Literature, source context and scope}
The primary AIM notes describe a brainstorming session. The original excerpt and all catalogue qualifications remain above. They do not specify a mathematical theorem or a complete empirical study. The formal model below is an explicitly limited research route introduced for this attempt, not a reconstruction of a missing original conjecture. The supplied catalogue's earlier research summary was inspected as context; its named framework is not adopted as an established solution. The elementary result below is not claimed novel, and no educational effectiveness data were collected.

\subsection{Attempt: choosing support packages under a fixed resource budget}
The source supplies no training packages or measured benefits. I tested a limited selection model: package $i$ has a known nonnegative additive benefit $v_i$ and positive integer cost $c_i$, it may be purchased at most once, and total cost cannot exceed integer budget $B$. The tempting route is to rank packages by benefit per unit cost.
\begin{proposition}
The simple ratio-greedy rule can achieve an arbitrarily small fraction of optimal benefit. A dynamic program computes the exact optimum in $O(nB)$ arithmetic/comparison operations for $n$ packages.
\end{proposition}
\begin{proof}
For any integer $B>2$, offer two packages: $(c,v)=(B,B)$ and $(1,2)$. The second has ratio two, larger than the first's ratio one, so ratio-greedy selects it. The remaining budget $B-1$ cannot buy the first package; greedy benefit is two, whereas the optimum buys the first for benefit $B$. The ratio $2/B$ tends to zero.

Let $D_i(b)$ be the largest benefit obtainable from the first $i$ packages within budget $b$. Initialize $D_0(b)=0$. If $c_i>b$, set $D_i(b)=D_{i-1}(b)$; otherwise set
\[
D_i(b)=\max\{D_{i-1}(b),\ v_i+D_{i-1}(b-c_i)\}.
\]
Every feasible selection either omits package $i$ or includes it and leaves a feasible selection of earlier packages at the reduced budget. Both alternatives in the recurrence are attainable, so induction on $i$ proves equality with the optimum. There are $(n+1)(B+1)$ states and constant work per transition.
\end{proof}
\subsection{Verification and remaining gaps}
At budget 100 the example gives greedy benefit two and optimal benefit 100, reproduced by the recurrence. The running time is pseudopolynomial because $B$ may have only $O(\log B)$ input bits. Additivity, known benefits, costs and indivisible purchases are modeling choices; real training and ongoing support can interact or affect different outcomes. The algorithm does not identify what support instructors need and does not validate the record's unverified reconstruction.

\section{References and source attribution}
American Institute of Mathematics, workshop discussion notes on open educational resources and assessment, \url{https://aimath.org/pastworkshops/oerassessmentproblems.pdf}. The 16-page primary PDF was opened on 11 October 2026 and its numbered brainstorming item was checked against the supplied record. The source supports the provenance and agenda, not the additional modeling assumptions or an empirical educational conclusion. Catalogue and contributor attribution are retained in the reproduced Sources section. The worked mathematics is proved here without reliance on an unverified external theorem.


\paragraph{Integration review.} Codex (GPT-6; exact serving revision not exposed) reviewed the written arguments and expanded layout and source-reference wrappers. Repository metadata and the production web renderer were checked using Python 3.13 and JavaScript with MathJax 3.2.2. This model-assisted review is not independent human proof verification; no PDF compilation or formal proof checking is asserted.

\section{Completion Estimate}
\noindent\textbf{COMPLETION: 0\%}
No substantiated new progress toward the entire original target is claimed. The worked results and failed route are delimited above; a known result, illustrative model, or auxiliary obstruction is not counted as original-target completion.

\section{Final status and attribution}
\textbf{Attempt status: unresolved.} This describes this research attempt, not a new ruling on the source catalogue's status or on external literature claims.
The mathematics and this self-audit were generated by OpenAI \texttt{gpt-6-astra}, reasoning \texttt{ultra}, on 11 October 2026. Codex (GPT-6; exact serving revision not exposed) orchestrated the batch and contributed formatting/repository review. No human mathematical authorship, independent human proof checking, or formal verification is claimed. The historical preparation paragraphs in the reproduced catalogue body concern that body only.

Contribution prompt A was applied to the complete definition and supplied source record, followed by the adversarial checks of prompt B. Local reproduction files are \texttt{batch008/checks.py}, \texttt{computations.json}, and the per-ID audit JSON. They were not public links at preparation time. Numerical checks are finite and do not prove unrestricted results. No LaTeX engine was found locally; compilation was not performed by this research agent.

\end{document}
